\section{Résumé}

La réalisation de ce projet a permis de valider une structure électronique fiable et performante. Bien que la phase de conception ait exigé un investissement en temps considérable, la capacité de résolution du chemin et la vitesse de parcours obtenus à la fin de la présentation finale ont démontré la justesse de nos choix techniques. Le robot a circulé avec peu d'incident matériel, prouvant la stabilité de son circuit et l'efficacité du suivi de ligne, rendu possible par l'emplacement optimisé des capteurs. L'unique incident matériel a été causé par le positionnement des pastilles de couleurs empêchant de détecter la fin du virage. Par conséquent, une amélioration possible serait d'utiliser le capteur RGB pour aider à savoir quand un virage est terminé. Autrement, l'intégration des périphériques a été un succès : le capteur RGB a détecté les pastilles avec précision, tandis que l'écran LCD fournissait un retour d'information immédiat sur les couleurs identifiées.                                                                                                                                  
Le partie algorithmique a constitué le défi majeur de notre réflexion. Initialement, nous avions privilégié une approche par force brute pour sa simplicité d'implémentation, capable de trouver le chemin optimal dans des environnements de taille réduite. Cependant, nous avons rapidement identifié une limite critique : avec une complexité en $n!$, l'algorithme devenait inefficace sur des grilles plus larges, comme celles de dimension 7, nécessitant environ 3 minutes de calcul sur un ordinateur. Face à ce constat et disposant de quatre jours supplémentaires, nous avons optimisé notre solution en développant un algorithme basé sur le parcours en largeur (BFS) couplé à une logique de Voyageur de Commerce (TSP). Bien que cette seconde approche ait montré des difficultés de gestion des demi-tours le jour de la soutenance, notre stratégie de redondance nous a permis de basculer avec succès sur la force brute pour accomplir le parcours.


\section{Perspectives}

La réussite de ce projet repose également sur une organisation humaine solide et une gestion d'équipe adaptative. Nous avons abordé le travail en répartissant les tâches sous forme de "couches", ce qui a permis à chacun de s'approprier des fonctions primaires tout en maintenant une vision globale du système. La partie électronique a été le fruit d'un effort collectif intense, où le circuit a été refait et débogué à plusieurs reprises par l'ensemble des membres du groupe. Cette synergie s'est retrouvée dans le développement du code, où les frontières entre les missions individuelles se sont effacées au profit de l'entraide.

En conclusion, nous avons su surmonter les difficultés inhérentes au travail en groupe en acceptant de "s'entre-couper" dans nos tâches respectives. Cette flexibilité nous a permis de repasser sur le travail des autres pour corriger, ajouter ou déboguer des éléments critiques, garantissant ainsi la cohérence finale du robot. Ce projet nous aura appris l'importance de la modularité, non seulement dans le code et l'électronique, mais aussi dans la collaboration humaine pour mener à bien un système complexe dans un temps limité.

\begin{table}[h]
\centering
\small % Réduit la taille de la police
\setlength{\tabcolsep}{3pt} % Réduit l'espace entre les colonnes (défaut à 6pt)
\renewcommand{\arraystretch}{1.5}
\begin{tabular}{|l|*{6}{c|}}
\hline
\textbf{Tâches} & \textbf{Noémie} & \textbf{Aleks} & \textbf{Adam} & \textbf{Youssef} & \textbf{Mathis} & \textbf{Bilal} \\ \hline
Construction du robot & X & X & X & X & X & X \\ \hline
Fonctions moteur & X &  &  & X &  &X \\ \hline
Fonctions capteurs & X &  & & X &  & X \\ \hline
Fonctions LCD & X & X &  &  &  & \\ \hline
Tests unitaires électronique & X & X & & & X & \\ \hline
Débogage électronique & X &  &  & X &  & X \\ \hline
Analyse Descendante & X & & & &  & X\\ \hline
Conception préliminaire &  & X & X & & X & X \\ \hline
Fonctions liées à la force brute &  & & & & X & \\ \hline
Fonctions liées au parcours en largeur & X & X & X & & &  \\ \hline
Tests unitaires algorithme &  & X  & & X & &  \\ \hline
Débogage algorithme &  & X & & & X &  \\ \hline
Mise en place du makefile et bashrc & & & & X & X &  \\ \hline
Rédaction du rapport (partie Electronique) & X &  &  &  &  &  \\ \hline
Rédaction du rapport (partie Informatique) & X & X & X & X & X & X \\ \hline
Documentation &  & X &  &  & X &  \\ \hline
\end{tabular}
\caption{Répartition des tâches}
\label{tab:repartition_taches}
\end{table}


\begin{thebibliography}{9}

\bibitem{rpi} 
Raspberry Pi Foundation, \textit{Power Consumption Benchmarks}. \\
\url{https://www.pidramble.com/wiki/benchmarks/power-consumption}

\bibitem{ttmotor} 
Adafruit Industries, \textit{DC Gearbox Motor "TT Motor" Datasheet}. \\
\url{https://media.digikey.com/pdf/Data%20Sheets/Adafruit%20PDFs/3777_Web.pdf}

\bibitem{l293d} 
Texas Instruments, \textit{L293D Quadruple Half-H Drivers Datasheet}. \\
\url{https://www.ti.com/lit/ds/symlink/l293.pdf}

\bibitem{lcd} 
Waveshare, \textit{LCD1602 I2C Module Parameters}. \\
\url{https://www.waveshare.com/wiki/LCD1602_I2C_Module}

\bibitem{rgb} 
OSRAM, \textit{TCS3472 Color Light-to-Digital Converter Datasheet}. \\
\url{https://cdn-shop.adafruit.com/datasheets/TCS34725.pdf}

\bibitem{tcrt} Vishay Semiconductors, \textit{TCRT5000 Reflective Optical Sensor Datasheet}. \\
\url{https://www.vishay.com/docs/83760/tcrt5000.pdf}

\end{thebibliography}
